\section{スーパーミラーによるスピンの選択}
\nocite{neutron_spin_optics}
この実験では、特定のスピンを持つ中性子を選択的に取り出す必要がある。
この節では、そのために用いるスーパーミラーの原理について説明する。

\subsection{中性子の光学的性質}
\paragraph{屈折率}
中性子が物質中で感じるポテンシャルを$V$とすると、エネルギー保存から
\begin{align}
k^2-k'^2=2mV
\end{align}
となる。$k$は入射中性子の波数、$k'$は物質中での中性子の波数である。
屈折率の定義
\begin{align}
n=\frac{k'}{k}
\end{align}
から、中性子の物質中における屈折率は
\begin{align}
n^2=1-\frac{2mV}{k^2}\label{mirror_neutron_refindex}
\end{align}
となる。

\paragraph{全反射が起きる条件}\label{mirrir_perfect_reflection}
$n-1$が有限の値を持つことは、中性子が全反射しうることを意味する。全反射が起きるための角度の条件は、
臨界角
\begin{align}
\theta^*=\arccos{n}
\end{align}
を用いて
\begin{align}
n\leq\cos\theta^* \label{mirror_refindex_range}
\end{align}
となる。また、全反射が起きるための入射エネルギー$E$の条件は、(\ref{mirror_neutron_refindex}), (\ref{mirror_refindex_range})より
\begin{align}
&n^2=1-\frac{2mV}{k^2}\leq\cos^2\theta^*\\
&E\sin^2\theta\leq E\sin^2\theta^*=\frac{k^2}{2m}\sin^2\theta^*\leq V
\end{align}
となる。すなわち、臨界角以下で入射する時、エネルギーの``ミラーに対し垂直な成分''が$V$よりも小さければ、
全反射が起きることがわかる。

\subsection{磁性体単層膜によるスピンの選択}
磁性体の単層膜を利用することで、特定のスピンを持つ中性子を選択的に取り出すことができる。
中性子が単層膜中で受けるポテンシャルは、核力によるポテンシャル$V_\mathrm{n}$、単層膜中の磁束密度$B$を用いて
\begin{align}
V^{\pm}&=V_\mathrm{n}\pm|\mu_\mathrm{n}|B
\end{align}
となる。$\mu_\mathrm{n}|B|$の符号とスピンの向きの対応は磁場の向きによって決まるが、ここでは上向きスピンのときに正になるものとする。
すなわち、上向きスピンの中性子は$V^+$、下向きスピンの中性子は$V^-$を感じる。

\ref{mirrir_perfect_reflection}節で述べたように、入射中性子のエネルギーを$E$とすると、
\[
E\sin^2\theta\leq V
\]
のときに全反射が起きる。
$V^-< E\sin^2\theta\leq V^+$のエネルギーを持つ中性子をこの単層膜に入射させると、下向きスピンの中性子はほとんどが透過するが、
上向きスピンの中性子は全反射される。$V^+<E\sin^2\theta$の中性子が入射した場合は、上向き、下向き両方の粒子が
透過する確率を持つ。
そのため、透過した中性子には上向きスピンと下向きスピンの両方が含まれる。
$E\sin^2\theta\leq V^-$となるような低エネルギーの中性子は今回の実験では無視できるほど少ないため、
考えなくて良い。

\begin{figure}[h]
\centering
\includegraphics[width=8cm]{mirror/mono_mirror.pdf}
\caption{単層膜によるスピンの選択の原理。実際には、ポテンシャルは紙面奥に向かって2次元に広がっており、
中性子は臨界角以下で入射していることに注意。}
\end{figure}
このようにして、単層膜は上向きスピンの粒子のみを選択的に反射する。

\subsection{スーパーミラー}
\begin{figure}[H]
\centering
\includegraphics{mirror/super_mirror.pdf}
\caption{スーパーミラーの構造\cite{seki}\label{mirror_super_mirror}}
\end{figure}
\begin{figure}[h]
\centering
\includegraphics[width=9cm]{mirror/super_mirror_potential.pdf}
\caption{スーパーミラーのポテンシャル。破線は磁場によらない核力によるポテンシャル。実線（上）は上向きスピンの中性子が感じるポテンシャル。
実線（下）は下向きスピンの中性子が感じるポテンシャル。空気やシリコンの透磁率は小さいため、全体に磁場をかけても
ポテンシャルはほとんど影響を受けない。}
\end{figure}
図\ref{mirror_super_mirror}のように膜厚を少しずつ変えた多層膜を使うことで、
単層膜よりも高いエネルギーの中性子を反射するミラーを作ることができる。
これをスーパーミラーと呼ぶ。

\paragraph{Bragg反射による偏極}
$E\sin^2\theta\leq V^+$の中性子は多層膜の表面で全反射されるが、$V^+<E\sin^2\theta$の中性子は多層膜の内部に侵入する。
多層膜の膜間隔を$d$とすると、侵入した中性子はBragg条件
\begin{align}
2d\sin\theta=\lambda
\end{align}
を満たす場合に反射波が強め合う。$\lambda$は中性子の物質波の波長で、
$\lambda=2\pi/k$の関係にある。

$d$を少しずつ変えることで、$V^+<E\sin^2\theta$の様々なエネルギーを持つ上向きスピンの中性子がBragg反射される。
下向きスピンの中性子もわずかに反射されるが、その数は少なく、無視できる。
こうすることで、単層膜に比べより広いエネルギーの中性子のスピンを偏極することができる。
